\honest{2-Place and 1-Place Words}{Eliezer Yudkowsky}{2008}

I have written before about the pulp covers on which a bug-eyed monster carries off a woman in a torn dress. Let us call the artist Fred. Fred says the monster ``can see that human females have soft, pleasant skin instead of slimy scales''; it may be an alien, but it is not stupid enough to make such a basic mistake about sexiness.

Fred's error is to treat sexiness, a function of two arguments, an admirer and an entity, as a function of the entity alone, so that it seems to depend on nothing else; this is the Mind Projection Fallacy. It is like asking whether a building is on the left side of the road without saying which way you are going. \nb{Philosophers state relativism in the same form: the Stanford Encyclopedia describes predicates such as ``right'' and ``beautiful'' as apparently one-place but ``elliptical for two-place predicates.'' Response-dependent theories of value index properties to subjects in the same way.}

An ``equally valid'' view is that each speaker's ``sexy'' names a different one-place function, which need not mention the speaker. Who says that Fred and Bloogah mean the same thing by ``sexy''? Fred's function can be specified in curves, skin, clothing and status cues, without mentioning Fred, and John might use the same one. Fred uses Sexiness\_20934, the monster Bloogah uses Sexiness\_72546, and anyone can check that the woman scores 5 on the first and 0.01 on the second, while a slime mold scores 3 on the second. Which function Fred uses is an empirical fact about Fred. \nb{For moral words this view is called contextualism. Its best-known difficulty is that two speakers using different functions do not contradict each other, even when they seem to disagree.}

Currying joins the two views. Give the two-place function plus the number 2 and you get a one-place function that adds 2; feed a seven-place function four arguments and you get a three-place one. A purist would say every function takes one argument: plus takes an integer and returns a function from integers to integers. Give ``sexiness'' an admirer and you get that admirer's fixed function, a mathematical object; that the admirer's intuitions compute it is an empirical fact about the admirer, and its value on the woman is a fact about her. \nb{This is the structure of David Kaplan's distinction between character, a function from contexts, and content.}

The same goes for Putnam's Twin Earth, where the watery stuff is XYZ, set centuries ago so that no one can test it. Some said ``water'' means the same on both planets, since people have the same sensory test in mind; others said it means H2O here and XYZ there. Read ``water'' as a function that eats a world, and it is the same concept on both planets; read it as the result for our world, and Twin Earth has no water. There is no point arguing over what the syllables really mean. Picking one definition is not enough to avoid confusion; you must train yourself to be aware of the curried and uncurried forms. ``The whole confusingness of the subsequent philosophical debate, rested on a tendency to instinctively curry concepts or instinctively uncurry them.'' I give no argument for this. \nb{Two-dimensional semantics had proposed this two-level answer, and it has been criticized on grounds about a priori knowledge and meaning, not only about instinct.}

Fred's brain simply computes his function and calls the result a property of the woman; to see that the monster chooses by a different function, Fred must consciously re-envision sexiness as two-place. The monster would compute Fred's function the same way, if it cared, but it uses a different one to decide whom to kidnap. We will need all this later, to taboo ``objective'', ``subjective'' and ``arbitrary''.
